\documentclass{article}
\usepackage{amsmath,amssymb,amsthm}
\usepackage{algorithm}
\usepackage{algpseudocode}
\usepackage{hyperref}
\usepackage{enumitem}
\newtheorem{theorem}{Theorem}
\newtheorem{lemma}{Lemma}
\newtheorem{corollary}{Corollary}
\newtheorem{assumption}{Assumption}
\newtheorem{definition}{Definition}
\newcommand{\E}{\mathbb{E}}
\newcommand{\R}{\mathbb{R}}
\newcommand{\inner}[2]{\langle #1,#2\rangle}
\newcommand{\argmin}{\operatorname*{arg\,min}}

\begin{document}

\section*{Frank–Wolfe Conditional Gradient Method}

\section*{Mathematical Formulation}
Let $f:\R^{d}\rightarrow \R$ be a convex and $L$‑smooth function and let $C\subset\R^{d}$ be a non‑empty compact convex set (the feasible domain).  
Given a current iterate $x_{k}\in C$, the Frank–Wolfe method performs

\[
\begin{aligned}
s_{k} &:= \argmin_{s\in C}\;\inner{\nabla f(x_{k})}{s} \quad\text{(Linear Minimization Oracle)}\tag{1}\\[4pt]
\gamma_{k} &\in[0,1]\quad\text{(stepsize, either prescribed or obtained by line search)}\\[4pt]
x_{k+1} &:= (1-\gamma_{k})\,x_{k}+\gamma_{k}\,s_{k}\qquad\in C.\tag{2}
\end{aligned}
\]

A common choice is the \emph{open‑loop} stepsize
\[
\gamma_{k}= \frac{2}{k+2},\qquad k=0,1,2,\dots
\]
or the exact line‑search
\[
\gamma_{k}= \argmin_{\gamma\in[0,1]} f\bigl((1-\gamma)x_{k}+\gamma s_{k}\bigr).
\]

\section*{Pseudocode}
\begin{algorithm}[H]
\caption{Frank–Wolfe Conditional Gradient}
\begin{algorithmic}[1]
\Require Convex $f$, compact convex $C$, initial $x_{0}\in C$, stepsize rule $\{\gamma_{k}\}$, max iterations $K$.
\For{$k=0$ \textbf{to} $K-1$}
    \State Compute gradient $g_{k}\gets\nabla f(x_{k})$.
    \State \textbf{LMO:} $s_{k}\gets\argmin_{s\in C}\inner{g_{k}}{s}$.
    \State Choose stepsize $\gamma_{k}\in[0,1]$ (e.g. $\gamma_{k}=2/(k+2)$ or line search).
    \State Update $x_{k+1}\gets (1-\gamma_{k})x_{k}+\gamma_{k}s_{k}$.
\EndFor
\State \Return $x_{K}$ (or the best iterate according to $f$).
\end{algorithmic}
\end{algorithm}

\section*{Dry Run Example}
Consider the one‑dimensional problem  

\[
f(w)=(w-3)^{2},\qquad C=[0,5],\qquad w_{0}=0,
\]
and use the open‑loop stepsize $\gamma_{k}=2/(k+2)$.  
Because $C$ is an interval, the linear minimization oracle reduces to checking the two endpoints.

\begin{center}
\begin{tabular}{c|c|c|c|c}
$k$ & $\nabla f(w_{k})$ & $s_{k}$ (LMO) & $\gamma_{k}=2/(k+2)$ & $w_{k+1}$ \\ \hline
0 & $2(w_{0}-3)=-6$ & $\argmin_{s\in[0,5]}(-6)s =0$ & $2/2=1$ & $(1-1)0+1\cdot0=0$ \\
1 & $-6$ (still) & $0$ & $2/3\approx0.667$ & $(1-0.667)0+0.667\cdot0=0$ \\
2 & $-6$ & $0$ & $2/4=0.5$ & $0$ 
\end{tabular}
\end{center}

In this trivial example the oracle always returns the left endpoint $0$, therefore the iterates stay at $0$ and $f(w_{k})=9$ for all $k$.  
If we start from $w_{0}=5$, the oracle would return $s_{k}=5$ for the first step (since $\nabla f(5)=4>0$) and the method would move towards the optimum $3$.

\newpage
\section*{Assumptions}
\begin{assumption}[Convexity]\label{ass:convex}
$f$ is convex on $\R^{d}$, i.e.
\[
f(y)\ge f(x)+\inner{\nabla f(x)}{y-x},\qquad\forall x,y\in\R^{d}.
\]
\end{assumption}
\begin{assumption}[Smoothness]\label{ass:smooth}
$f$ is $L$‑smooth on $C$, i.e. for all $x,y\in C$,
\[
f(y)\le f(x)+\inner{\nabla f(x)}{y-x}+\frac{L}{2}\|y-x\|^{2}.
\]
\end{assumption}
\begin{assumption}[Compact Feasible Set]\label{ass:compact}
$C$ is compact and convex with diameter
\[
D:=\max_{x,y\in C}\|x-y\|<\infty .
\]
\end{assumption}
\begin{definition}[Curvature Constant]\label{def:curv}
The curvature of $f$ over $C$ is
\[
C_{f}:=\sup_{\substack{x,s\in C\\ \gamma\in[0,1]}}
\frac{2}{\gamma^{2}}\Bigl(f\bigl(x+\gamma(s-x)\bigr)-f(x)-\gamma\inner{\nabla f(x)}{s-x}\Bigr).
\]
If $f$ is $L$‑smooth on $C$, then $C_{f}\le L D^{2}$.
\end{definition}

\section*{Main Convergence Theorem}
\begin{theorem}[Primal Gap Decay]\label{thm:fw}
Let Assumptions \ref{ass:convex}–\ref{ass:compact} hold and let $\{x_{k}\}$ be generated by the Frank–Wolfe method with the open‑loop stepsize $\gamma_{k}=2/(k+2)$.  
Denote $x^{*}\in\arg\min_{x\in C}f(x)$. Then for all $k\ge0$,
\[
f(x_{k})-f(x^{*})\;\le\;\frac{2C_{f}}{k+2}.
\tag{3}
\]
Consequently $f(x_{k})\to f(x^{*})$ at a rate $\mathcal{O}(1/k)$.
\end{theorem}

\section*{Theoretical Properties}
\begin{itemize}[leftmargin=*]
\item \textbf{Stability.} The iterates always remain in $C$ because $x_{k+1}$ is a convex combination of two points in $C$.
\item \textbf{Stepsize Sensitivity.} The $\mathcal{O}(1/k)$ rate holds for any stepsize sequence satisfying $\gamma_{k}\in[0,1]$, $\sum_{k}\gamma_{k}=+\infty$, and $\sum_{k}\gamma_{k}^{2}<\infty$. The specific choice $\gamma_{k}=2/(k+2)$ yields the tight constant $2C_{f}$.
\item \textbf{Comparison to Projected Gradient Descent.} While projected gradient needs an orthogonal projection onto $C$, Frank–Wolfe only requires the linear minimization oracle, which is often cheaper for polytopes or atomic sets. The price is a slower $\mathcal{O}(1/k)$ rate compared with $\mathcal{O}(1/k)$ for projected gradient under the same smoothness, but Frank–Wolfe enjoys sparsity of iterates.
\item \textbf{Special Cases.} If $f$ is also strongly convex with parameter $\mu>0$, the primal gap can be improved to $\mathcal{O}(1/k^{2})$ under a variant with away steps; however the basic Frank–Wolfe algorithm retains the $\mathcal{O}(1/k)$ bound.
\end{itemize}

\section*{Discussion}
The theorem shows that the Frank–Wolfe algorithm converges sublinearly without ever leaving the feasible set. The curvature constant $C_{f}$ captures both the smoothness of $f$ and the geometry of $C$. When $C$ is a simplex or a spectahedron, the LMO is extremely cheap, making the method attractive for large‑scale sparse optimization.

\newpage
\section*{Preliminaries (Definitions and Lemmas)}
\begin{lemma}[Smoothness Inequality]\label{lem:smooth}
If $f$ is $L$‑smooth on $C$, then for any $x,y\in C$,
\[
f(y)\le f(x)+\inner{\nabla f(x)}{y-x}+\frac{L}{2}\|y-x\|^{2}.
\]
\end{lemma}
\begin{lemma}[Descent Lemma via Curvature]\label{lem:curv}
For any $x\in C$, $s\in C$ and $\gamma\in[0,1]$,
\[
f\bigl(x+\gamma(s-x)\bigr)\le f(x)+\gamma\inner{\nabla f(x)}{s-x}+\frac{\gamma^{2}}{2}C_{f}.
\tag{4}
\]
\end{lemma}
\begin{proof}
By definition of $C_{f}$,
\[
C_{f}\ge\frac{2}{\gamma^{2}}\bigl(f(x+\gamma(s-x))-f(x)-\gamma\inner{\nabla f(x)}{s-x}\bigr),
\]
which rearranges to (4). \hfill $\square$
\end{proof}

\section*{Main Proof (Step‑by‑Step Derivation)}
We prove Theorem \ref{thm:fw}.  Throughout the proof we use the notation $g_{k}:=\nabla f(x_{k})$.

\begin{enumerate}
\item[Step 1.] \textbf{Optimality of the LMO.}  
Because $s_{k}$ solves \eqref{1},
\[
\inner{g_{k}}{s_{k}} \le \inner{g_{k}}{x^{*}},\qquad\forall x^{*}\in C.
\tag{5}
\]

\item[Step 2.] \textbf{One‑step progress via curvature.}  
Apply Lemma \ref{lem:curv} with $x=x_{k}$, $s=s_{k}$ and $\gamma=\gamma_{k}$:
\[
f(x_{k+1})\le f(x_{k})+\gamma_{k}\inner{g_{k}}{s_{k}-x_{k}}+\frac{\gamma_{k}^{2}}{2}C_{f}.
\tag{6}
\]

\item[Step 3.] \textbf{Replace the linear term using (5).}  
From (5) we have $\inner{g_{k}}{s_{k}-x_{k}}\le\inner{g_{k}}{x^{*}-x_{k}}$. Substituting into (6) yields
\[
f(x_{k+1})\le f(x_{k})+\gamma_{k}\inner{g_{k}}{x^{*}-x_{k}}+\frac{\gamma_{k}^{2}}{2}C_{f}.
\tag{7}
\]

\item[Step 4.] \textbf{Use convexity of $f$.}  
Convexity (Assumption \ref{ass:convex}) gives
\[
f(x^{*})\ge f(x_{k})+\inner{g_{k}}{x^{*}-x_{k}}.
\tag{8}
\]
Rearranging (8) we obtain $\inner{g_{k}}{x^{*}-x_{k}}\le f(x^{*})-f(x_{k})$.

\item[Step 5.] \textbf{Combine (7) and (8).}  
Insert the bound from (8) into (7):
\[
f(x_{k+1})\le f(x_{k})+\gamma_{k}\bigl(f(x^{*})-f(x_{k})\bigr)+\frac{\gamma_{k}^{2}}{2}C_{f}.
\tag{9}
\]

\item[Step 6.] \textbf{Define the primal gap.}  
Let $\Delta_{k}:=f(x_{k})-f(x^{*})\ge0$.  
Equation (9) becomes
\[
\Delta_{k+1}\le (1-\gamma_{k})\Delta_{k}+\frac{\gamma_{k}^{2}}{2}C_{f}.
\tag{10}
\]

\item[Step 7.] \textbf{Insert the open‑loop stepsize.}  
Set $\gamma_{k}=2/(k+2)$. Then $1-\gamma_{k}= \frac{k}{k+2}$ and $\gamma_{k}^{2}=4/(k+2)^{2}$. Hence (10) reads
\[
\Delta_{k+1}\le\frac{k}{k+2}\Delta_{k}+\frac{2}{(k+2)^{2}}C_{f}.
\tag{11}
\]

\item[Step 8.] \textbf{Induction hypothesis.}  
We claim that for all $k\ge0$,
\[
\Delta_{k}\le\frac{2C_{f}}{k+2}.
\tag{12}
\]

\item[Step 9.] \textbf{Base case $k=0$.}  
From definition $\Delta_{0}=f(x_{0})-f(x^{*})\le \frac{2C_{f}}{2}=C_{f}$, which is true because $C_{f}\ge0$ and the curvature definition ensures $f(x_{0})-f(x^{*})\le C_{f}$ (a direct consequence of Lemma \ref{lem:curv} with $\gamma=1$). Thus (12) holds for $k=0$.

\item[Step 10.] \textbf{Inductive step.}  
Assume (12) holds for some $k$. Plug into (11):
\[
\Delta_{k+1}\le\frac{k}{k+2}\cdot\frac{2C_{f}}{k+2}+\frac{2C_{f}}{(k+2)^{2}}
= \frac{2C_{f}}{(k+2)^{2}}\bigl(k+1\bigr)
= \frac{2C_{f}}{k+3}.
\tag{13}
\]
Since $k+3 = (k+1)+2$, (13) is exactly the statement (12) for index $k+1$. Hence the induction is complete.

\item[Step 11.] \textbf{Conclusion.}  
Inequality (12) is equivalent to the bound (3) of Theorem \ref{thm:fw}. Thus the primal gap decays as $\mathcal{O}(1/k)$.
\end{enumerate}
All steps rely solely on Assumptions \ref{ass:convex}–\ref{ass:compact} and on the definition of the curvature constant.

\section*{Rate Derivation (With Constants)}
From (12) we have explicitly
\[
f(x_{k})-f(x^{*})\le\frac{2C_{f}}{k+2}\le\frac{2L D^{2}}{k+2},
\]
where we used $C_{f}\le L D^{2}$ (see Definition \ref{def:curv}). Hence the constant in the $\mathcal{O}(1/k)$ rate can be taken as $2L D^{2}$.

\section*{Special Cases}
\begin{enumerate}[label=\arabic*.]
\item \textbf{Exact line search.}  
If $\gamma_{k}$ is chosen by minimizing $f\bigl((1-\gamma)x_{k}+\gamma s_{k}\bigr)$ over $\gamma\in[0,1]$, the same recurrence (10) holds with $\gamma_{k}$ replaced by the optimal $\gamma_{k}^{\star}$, and the bound (3) remains valid (often with a smaller constant).

\item \textbf{Away‑step Frank–Wolfe.}  
When $C$ is a polytope, adding away steps yields a linear convergence rate $\mathcal{O}((1-\rho)^{k})$ for some $\rho>0$ under a geometric condition (the so‑called pyramidal width). The basic proof follows the same template but augments the curvature analysis with a “drop step”.

\item \textbf{Strongly convex $f$.}  
If $f$ is $\mu$‑strongly convex, one can combine the curvature bound with strong convexity to obtain
\[
\Delta_{k}\le\frac{2C_{f}}{\mu(k+2)},
\]
which still is $\mathcal{O}(1/k)$ but the constant improves proportionally to $1/\mu$.
\end{enumerate}

\end{document}